\begin{ouverture}[calcul]{10}{Inconnues et équations simples}{Partie 2 · Algèbre et grandeurs}
{Opération à trou · Addition · Multiplication · Vérification}
\objectif{Donner à la lettre le statut d’inconnue.}
\objectif{Résoudre $x+b=c$ ou $ax=c$ par opérations inverses.}
\objectif{Traduire un problème simple et contrôler la solution.}
\prerequis{Nombres relatifs et quotients.}
\end{ouverture}
\pagestyle{otMain}\essentiel\fontsize{11}{15}\selectfont
\phantomsection\label{cours-10-1}
\essSection{1}{Retrouver le nombre caché}
\begin{definition}Une inconnue est un nombre à déterminer. Une équation est une égalité où intervient une inconnue. Une solution est une valeur qui rend l’égalité vraie. Résoudre signifie trouver toutes les solutions demandées.\end{definition}
\label{fig-cours-10-1-cours-equation}\figureGeo{cours-equation}

\begin{savoirfaire}[Utiliser les opérations inverses]Pour $x+7=19$, retirer 7 de 19 : $x=12$. Pour $5x=18$, diviser 18 par 5 : $x=18/5=3{,}6$. Pour $\frac x3=1$, c’est-à-dire $x:3=1$, multiplier par 3 : $x=3$. Ce sont les opérations inverses de l’addition, de la multiplication et de la division.\end{savoirfaire}
\begin{attention}$5x$ signifie cinq fois $x$, pas $5+x$. Une valeur testée qui ne fonctionne pas n’est pas une solution. Cas particuliers : $2x=0$ et $\frac x2=0$ ont pour seule solution 0, car $0:2=0$ et $0\times2=0$.\end{attention}
\clearpage
\phantomsection\label{cours-10-2}
\essSection{2}{Modéliser et contrôler}
\begin{exemple}Trois carnets identiques coûtent 10,50 euros. Si $x$ est le prix d’un carnet, $3x=10{,}50$, donc $x=3{,}50$. Le contrôle $3\times3{,}50=10{,}50$ confirme le résultat et son unité.\end{exemple}
\begin{savoirfaire}[Quatre étapes]Nommer l’inconnue avec son unité. Traduire la relation connue par une égalité. Utiliser l’opération inverse. Remplacer la lettre par la réponse dans l’égalité de départ, puis répondre avec une phrase.\end{savoirfaire}
\begin{exemple}Pour $x+(-4)=3$, $x=3-(-4)=7$. Pour $x+6=2$, $x=2-6=-4$. La solution peut être négative si le contexte l’autorise.\end{exemple}
\begin{attention}Ce chapitre se limite aux équations du type $x+b=c$ ou $ax=c$, avec un coefficient $a$ non nul, décimal ou fraction simple : $\frac x3$ est le tiers de $x$. Les problèmes à étapes se résolvent aussi par raisonnement arithmétique, en remontant les opérations.\end{attention}

\clearpage\rubrique{Découverte}
\begin{blocLivre}{ex-10-1}\begin{exercice}[1]{Une addition à trou}
Résoudre $x+8=21$.
\end{exercice}\end{blocLivre}
\begin{blocLivre}{ex-10-2}\begin{exercice}[1]{Une multiplication à trou}
Résoudre $6x=42$.
\end{exercice}\end{blocLivre}
\begin{blocLivre}{ex-10-3}\begin{exercice}[1]{Un quotient exact}
Résoudre $4x=7$ en donnant une valeur exacte.
\end{exercice}\end{blocLivre}
\begin{blocLivre}{ex-10-4}\begin{exercice}[1]{Tester un candidat}
5 est-il une solution de $3x=18$ ?
\end{exercice}\end{blocLivre}
\begin{blocLivre}{ex-10-5}\begin{exercice}[1]{Une valeur négative}
Résoudre $x+9=4$.
\end{exercice}\end{blocLivre}

\rubrique{Entraînement}
\begin{blocLivre}{ex-10-6}\begin{exercice}[2]{Avec un décimal}
Résoudre $x+2{,}7=8{,}15$.
\end{exercice}\end{blocLivre}
\begin{blocLivre}{ex-10-7}\begin{exercice}[2]{Un coefficient décimal}
Résoudre $0{,}8x=6$.
\end{exercice}\end{blocLivre}
\begin{blocLivre}{ex-10-8}\begin{exercice}[2]{Additionner un négatif}
Résoudre $x+(-6)=2$.
\end{exercice}\end{blocLivre}
\begin{blocLivre}{ex-10-9}\begin{exercice}[2]{Trois équations proches}
Résoudre $7+x=12$, $7x=12$ et $\frac x7=12$.
\end{exercice}\end{blocLivre}
\begin{blocLivre}{ex-10-10}\begin{exercice}[2]{Une mauvaise opération}
Noé résout $5x=23$ en écrivant $x=23-5=18$. Corriger.
\end{exercice}\end{blocLivre}

\rubrique{Problèmes}
\begin{blocLivre}{ex-10-11}\begin{exercice}[3]{Un prix unitaire}
Neuf cahiers identiques coûtent 25,20 euros. Déterminer le prix d’un cahier en posant une équation.
\end{exercice}\end{blocLivre}
\begin{blocLivre}{ex-10-12}\begin{exercice}[3]{Une quantité initiale}
Après avoir ajouté 17 L, une cuve contient 63 L. Combien contenait-elle avant ?
\end{exercice}\end{blocLivre}
\begin{blocLivre}{ex-10-13}\begin{exercice}[3]{Aire et largeur}
Un rectangle de longueur 8 cm a une aire de 46 cm$^2$. Déterminer sa largeur.
\label{fig-ex-10-13}\figureGeo{eq-13}
\end{exercice}\end{blocLivre}
\begin{blocLivre}{ex-10-14}\begin{exercice}[3]{Parts dans un schéma}
Trois parts identiques et un supplément de 6 euros forment 45 euros. Trouver une part en raisonnant avec le schéma en barres.
\label{fig-ex-10-14}\figureGeo{eq-14}
\end{exercice}\end{blocLivre}
\begin{blocLivre}{ex-10-15}\begin{exercice}[3]{Une contrainte d’objets entiers}
On veut partager 50 jetons également entre huit personnes, sans casser de jeton. L’équation $8x=50$ a-t-elle une solution numérique ? Le partage sans reste est-il possible ?
\end{exercice}\end{blocLivre}

\clearpage
